\section{Fundamentação Teórica}
\label{sec:fundamentacao}

A fundamentação teórica deste trabalho apresenta os conceitos necessários para compreender o funcionamento de Algoritmos Genéticos (AGs), sua aplicação ao problema das \textit{n}-rainhas e o papel das estratégias de controle populacional na dinâmica evolutiva. Inicialmente, são discutidos os princípios básicos dos AGs, com ênfase em seus componentes estruturais. Em seguida, descreve-se o fluxo geral de um AG clássico, por meio de um pseudocódigo que servirá de referência conceitual para todas as variantes analisadas. Depois, revisita-se o problema das \textit{n}-rainhas como cenário experimental amplamente utilizado na literatura. São também apresentados os principais elementos do ambiente computacional adotado, com foco em bibliotecas numéricas para experimentação reprodutível. Por fim, são detalhadas as estratégias de controle populacional e os trabalhos relacionados que motivam o recorte específico deste estudo.

\subsection{Algoritmos Genéticos}

Algoritmos Genéticos (AGs) são meta-heurísticas inspiradas em princípios de evolução biológica, nas quais uma população de soluções candidatas é iterativamente modificada por operadores de seleção, cruzamento e mutação \cite{eiben2015introduction, holland1975adaptation}. Cada indivíduo é avaliado por uma função de aptidão (\textit{fitness}), e o processo evolutivo tende a favorecer configurações de melhor desempenho ao longo das gerações, compondo um ciclo de variação e seleção que equilibra exploração e intensificação no espaço de busca.

O comportamento de um AG é fortemente influenciado por três grupos de decisões de projeto \cite{lacerda1999introducao}:  
(i) a forma de representação dos indivíduos (ou codificação das soluções),  
(ii) a função de avaliação (função de aptidão) e  
(iii) os operadores evolutivos e a política de substituição populacional.  

A literatura tradicional em AGs dedica grande atenção à escolha da representação (binária, inteira, permutacional, real, entre outras) e ao desenho de operadores de cruzamento e mutação adequados ao problema \cite{goldberg1989genetic,eiben2015introduction}. Em contraste, o \textit{controle populacional} — isto é, a forma como indivíduos são mantidos, substituídos, removidos ou preservados ao longo das gerações — costuma ser tratado como um detalhe de implementação, apesar de influenciar diretamente aspectos como diversidade, pressão seletiva, convergência e risco de aprisionamento em ótimos locais \cite{whitley1994tutorial,mitchell1998introduction}.  

Neste trabalho, representação, função de aptidão, seleção, cruzamento e mutação são mantidos fixos, e o foco recai exclusivamente sobre o componente de controle populacional. Essa escolha permite isolar e mensurar o impacto de diferentes políticas de substituição no desempenho global do algoritmo, mantendo constantes os demais fatores que tradicionalmente afetam a qualidade da busca.

\subsection{Fluxo Geral de um Algoritmo Genético}

Embora existam inúmeras variantes, o fluxo de execução de um AG tradicional segue uma estrutura geral bem estabelecida, composta por inicialização da população, avaliação, reprodução e atualização da população ao longo de múltiplas gerações \cite{eiben2015introduction,mitchell1998introduction}. Critérios de parada usuais incluem: (i) número máximo de gerações, (ii) orçamento máximo de avaliações da função de aptidão e (iii) limite de tempo de execução. 

O pseudocódigo apresentado a seguir resume o processo e serve como modelo conceitual sobre o qual as diferentes estratégias de controle populacional serão discutidas. No contexto deste trabalho, o critério de parada é posteriormente especializado para um limite máximo de avaliações (ver Seção~\ref{sec:metodologia}).

\Needspace{20\baselineskip}

\begin{algorithm}[H]
\caption{Pseudocódigo de um Algoritmo Genético Clássico}
\label{alg:ag-classico}
\begin{algorithmic}[1]
\Statex \textbf{Entrada:} tamanho populacional $N$, orçamento máximo de avaliações $A_{\max}$, operadores genéticos
\Statex \textbf{Saída:} melhor solução encontrada
\Statex

\State Inicializar a população $P$ com $N$ indivíduos
\State Avaliar $P$ e inicializar o contador de avaliações $A \leftarrow N$
\While{critério de parada \textbf{não} satisfeito \textbf{e} $A < A_{\max}$}
    \Statex \textbf{// Seleção e Reprodução}
    \State Selecionar pais utilizando torneio (ou outro método de seleção)
    \State Aplicar operador de cruzamento para gerar descendentes
    \State Aplicar operador de mutação aos descendentes
    \State Avaliar os descendentes e atualizar $A$

    \Statex \textbf{// Controle Populacional}
    \State Combinar populações conforme a estratégia escolhida:
    \Statex \quad -- substituição total, ou
    \Statex \quad -- substituição parcial, ou
    \Statex \quad -- elitismo, ou
    \Statex \quad -- controle por idade (GAVaPS/Michalewicz)

    \State Atualizar a população $P$
\EndWhile
\State Retornar o melhor indivíduo de $P$
\end{algorithmic}
\end{algorithm}

Esse pseudocódigo é utilizado como referência conceitual para as variações testadas neste estudo, alterando-se apenas a estratégia de controle populacional aplicada em cada experimento, enquanto os demais componentes permanecem inalterados. A especialização para um critério de parada baseado exclusivamente em avaliações é detalhada na Seção~\ref{sec:metodologia}.

\subsection{O Problema das \textit{n}-Rainhas}

O problema das \textit{n}-rainhas consiste em posicionar \textit{n} rainhas em um tabuleiro de xadrez de dimensões \textit{n}×\textit{n} de forma que nenhuma delas ataque outra, o que implica ausência de conflitos em linhas, colunas e diagonais \cite{bell2009survey}. Embora a formulação seja simples, o número de configurações possíveis cresce rapidamente com \textit{n}, tornando abordagens exaustivas inviáveis para instâncias maiores e caracterizando um problema combinatório de difícil exploração.

Estudos como o de Ali et al. \cite{ali2023performance} mostram que buscas clássicas (BFS/DFS) são eficientes apenas para instâncias pequenas, ao passo que, para valores maiores de \textit{n}, meta-heurísticas, como AGs, tornam-se alternativas mais competitivas. Essa combinação de estrutura simples, dificuldade ajustável e ampla presença na literatura torna o problema das \textit{n}-rainhas um \textit{benchmark} natural para avaliar variações de AGs \cite{sarkar2017adaptive,farhan2015solving}.  

Neste trabalho, o problema das \textit{n}-rainhas é adotado como cenário experimental para comparar estratégias de controle populacional, beneficiando-se de uma representação permutacional que elimina conflitos de linha e coluna e concentra a função de aptidão na contagem de conflitos diagonais. Essa característica contribui para reduzir o impacto de soluções inviáveis e realçar o efeito das diferentes políticas de substituição populacional.

\begin{figure}[H]
  \centering
  \chessboard[
    setpieces={Qa1,Qb5,Qc8,Qd6,Qe3,Qf7,Qg2,Qh4},
    showmover=false,
    label=true
  ]
  \caption{Exemplo de solução válida para o problema das 8-rainhas.}
  \label{fig:tabuleiro-nrainhas}
\end{figure}

\subsection{Uso da Linguagem Python e Bibliotecas Relacionadas}

A implementação experimental foi desenvolvida em \texttt{Python}, linguagem amplamente adotada em computação científica pela clareza sintática e pelo ecossistema consolidado de bibliotecas numéricas \cite{python3reference}. Essa escolha facilita a construção de protótipos, a automatização de experimentos e a reprodutibilidade dos resultados.

A biblioteca \texttt{NumPy} foi utilizada para operações vetorizadas e para a representação eficiente da população \cite{harris2020array}. A geração de números aleatórios baseou-se em \texttt{numpy.random} \cite{python3random}, garantindo controle explícito de sementes e replicação de cenários. A visualização dos resultados foi realizada com \texttt{matplotlib} \cite{hunter2007matplotlib}, permitindo a construção de gráficos para análise comparativa das estratégias de controle populacional.

Esse conjunto de ferramentas permite concentrar o esforço no desenho dos experimentos e na análise do comportamento dos AGs, reduzindo interferências de infraestrutura e favorecendo a reprodutibilidade.

\subsection{Técnicas de Controle Populacional}
\label{sec:controlepopulacional}

\subsubsection{Substituição Total}

Na substituição total, a cada geração todos os indivíduos da população corrente são substituídos pelos descendentes recém-gerados. Em termos operacionais, o algoritmo:

\begin{enumerate}
    \item seleciona pares de pais na população atual;
    \item aplica cruzamento e mutação para produzir uma nova população de tamanho $N$;
    \item descarta integralmente a população anterior.
\end{enumerate}

Trata-se do esquema mais simples e amplamente utilizado em AGs, pois facilita a implementação e o controle de parâmetros \cite{goldberg1989genetic,whitley1994tutorial}. No entanto, por descartar toda a geração anterior, tende a reduzir a diversidade rapidamente, tornando o algoritmo mais suscetível à convergência prematura, especialmente quando a pressão seletiva é alta ou a mutação é conservadora \cite{dejong1975analysis,spears1993crossover}. A substituição total é frequentemente utilizada como referência de comparação para outras políticas de controle populacional.

\subsubsection{Substituição Parcial}

Na substituição parcial, apenas uma fração da população é renovada a cada geração, caracterizando um modelo do tipo \textit{steady-state} \cite{eiben2015introduction,whitley1994tutorial}. De forma geral, o processo segue os passos:

\begin{enumerate}
    \item a partir da população atual, selecionam-se pais e produzem-se alguns descendentes;
    \item esses descendentes são inseridos na população;
    \item um subconjunto de indivíduos é removido, preservando o tamanho populacional.
\end{enumerate}

A fração de indivíduos substituída em cada geração controla o ritmo de renovação: valores menores favorecem uma evolução mais gradual, enquanto valores maiores aproximam o comportamento do modelo geracional \cite{goldberg1989genetic,mitchell1998introduction}.  

Na literatura, diferentes critérios podem ser utilizados para decidir quais indivíduos permanecem ou são removidos (por aptidão, aleatoriedade, idade, entre outros). Neste trabalho, a estratégia de substituição parcial preserva uma fração fixa de indivíduos escolhidos \textbf{aleatoriamente}, sem preferência explícita pelos mais aptos. O elitismo, por sua vez, é tratado como uma estratégia distinta, responsável por garantir a preservação sistemática dos melhores indivíduos.

\subsubsection{Elitismo}

O elitismo introduz um mecanismo explícito de preservação dos melhores indivíduos \cite{goldberg1989genetic,whitley1994tutorial}. Em cada geração:

\begin{enumerate}
    \item uma pequena fração dos indivíduos com melhor aptidão é copiada diretamente para a próxima geração;
    \item o restante da população é preenchido com descendentes produzidos pelos operadores de cruzamento e mutação;
    \item indivíduos menos aptos são descartados para manter o tamanho populacional.
\end{enumerate}

Esse mecanismo garante que as melhores soluções encontradas não sejam perdidas por efeito do acaso, o que tende a acelerar a convergência \cite{dejong1975analysis,goldberg1989genetic}. Por outro lado, a preservação sistemática de indivíduos muito aptos pode reduzir a diversidade e acentuar o risco de estagnação em ótimos locais, sobretudo quando a população é pequena \cite{eiben2015introduction,lacerda1999introducao}. Assim, o elitismo representa um reforço da pressão seletiva que pode ser benéfico ou prejudicial, dependendo do equilíbrio com outros componentes do algoritmo.

\subsubsection{Controle Populacional por Idade (GAVaPS / Michalewicz)}

A estratégia de controle por idade baseia-se no modelo \textit{Genetic Algorithms with Varying Population Size} (GAVaPS) proposto por Michalewicz \cite{michalewicz1993genetic}. Nesse esquema, cada indivíduo passa a carregar, além do cromossomo e da aptidão, duas variáveis adicionais:

\begin{itemize}
\item sua \textbf{idade} (número de gerações sobrevividas);
\item um \textbf{tempo de vida} máximo (\textit{lifetime}), definido no momento do nascimento.
\end{itemize}

A premissa fundamental é que indivíduos mais aptos recebam um \textit{lifetime} maior, permanecendo por mais gerações na população, enquanto indivíduos menos aptos são naturalmente descartados mais cedo \cite{michalewicz1993genetic}. O cálculo do \textit{lifetime} é tipicamente obtido a partir da aptidão relativa do indivíduo, considerando estatísticas globais da população (mínimo, máximo e média). Essa aptidão é então mapeada por uma função — frequentemente bi-linear — para o intervalo discreto $[L_{\min}, L_{\max}]$. Assim, indivíduos de melhor desempenho recebem tempos de vida próximos a $L_{\max}$, ao passo que indivíduos fracos aproximam-se de $L_{\min}$.

O ciclo evolutivo do GAVaPS segue, de forma simplificada, três etapas principais \cite{michalewicz1993genetic,eiben2015introduction}:

\begin{enumerate}
\item incrementar a idade de todos os indivíduos;
\item remover aqueles cuja idade ultrapassou seu \textit{lifetime};
\item gerar novos descendentes — via seleção, cruzamento e mutação — em quantidade proporcional ao tamanho da população atual, permitindo que o número total de indivíduos varie dinamicamente.
\end{enumerate}

Diferentemente das estratégias tradicionais de substituição, em que a remoção é guiada pela aptidão instantânea, aqui a eliminação ocorre por envelhecimento. Mesmo indivíduos muito aptos eventualmente saem da população ao atingirem o limite de vida. Esse mecanismo introduz uma forma de \textit{diversidade temporal}, impedindo que poucas soluções dominem indefinidamente a população e promovendo um equilíbrio mais natural entre intensificação (preservar bons indivíduos por mais tempo) e exploração (renovar continuamente a população) \cite{michalewicz1993genetic}.

No contexto deste trabalho, o controle por idade é comparado diretamente com as políticas de substituição total, substituição parcial e elitismo, de modo a avaliar como a introdução do \textit{lifetime} e da variação dinâmica do tamanho populacional impacta a taxa de sucesso, o esforço computacional (em número de avaliações) e o comportamento evolutivo de AGs aplicados ao problema das \textit{n}-rainhas \cite{richter2012algoritmo,sarkar2017adaptive,ali2023performance}.
\subsection{Trabalhos Relacionados}

Diversos estudos utilizam o problema das \textit{n}-rainhas como \textit{benchmark} para investigar variações em AGs, especialmente no que diz respeito à representação e ao controle de diversidade. Esta subseção destaca trabalhos que fundamentam o desenho experimental adotado e situam a contribuição deste estudo, focado no impacto de diferentes estratégias de controle populacional.

\subsubsection{Representações Binária e Inteira para o Problema das \textit{n}-Rainhas \cite{richter2012algoritmo}}

Richter et al.\ comparam uma codificação binária e uma codificação inteira do tipo \textit{InOrder}. A primeira exige conversão e permite conflitos em linhas; a segunda mapeia diretamente colunas para linhas, evitando soluções inviáveis. Os autores mostram que a representação \textit{InOrder} reduz conflitos e acelera a convergência, reforçando a importância de uma representação adequada antes de comparar variações de um AG. O presente trabalho adota uma representação por permutação semelhante justamente para isolar o efeito das estratégias de substituição populacional.

\subsubsection{Comparação entre Algoritmo Genético e Buscas Clássicas (BFS/DFS) \cite{ali2023performance}}

Ali et al.\ comparam AG, BFS e DFS no problema das \textit{n}-rainhas. Para instâncias pequenas, as buscas clássicas são mais rápidas, mas o AG se torna mais competitivo à medida que o tamanho do tabuleiro cresce. Embora o foco seja a comparação entre paradigmas, o estudo reforça a adequação dos AGs para instâncias maiores. Diferentemente dos autores, este trabalho concentra-se exclusivamente em comparar estratégias dentro do próprio AG.

\subsubsection{Algoritmo Genético Adaptativo com Controle de Diversidade \cite{sarkar2017adaptive}}

Sarkar e Nag apresentam um AG adaptativo com mecanismos explícitos de preservação de diversidade, incluindo mutações mais frequentes e estruturas auxiliares para evitar estagnação. Os resultados mostram boa escalabilidade para $N$ elevados, sugerindo que o modo de renovação populacional influencia diretamente a convergência. Contudo, por combinar várias modificações simultâneas, o efeito de cada componente é difícil de isolar. Aqui, apenas a política de substituição é variada, mantendo-se constantes representação e operadores.

\subsubsection{Busca de Todas as Soluções Válidas para as 8-Rainhas \cite{farhan2015solving}}

Farhan et al.\ utilizam um AG para encontrar todas as 92 soluções da instância de 8 rainhas, empregando representação vetorial sem conflitos em linhas/colunas e avaliando apenas ataques diagonais. Embora o objetivo seja diferente — cobertura completa e não comparação entre estratégias — o trabalho ilustra a capacidade dos AGs de explorar múltiplas soluções válidas. Neste estudo, essa característica é explorada para observar como diferentes políticas de substituição e controle por idade afetam a taxa de sucesso ao encontrar ao menos uma solução válida.